\begin{abstract}
\noindent
An FFT-based spectrum analyzer can be inexpensive on average yet concentrate
most of its work in one audio-processing call. This paper describes how to
spread that work across successive sample calls without a separate analysis
thread. A resumable radix-2 FFT preserves its traversal state between calls;
windowing, spectrum reconstruction, smoothing, and output preparation are
scheduled alongside it. Each frame completes over one chosen update interval,
with a bounded operation count per call and an explicit delay before the
spectrum becomes available. We implement the method in two VCV Rack
visualizers and evaluate it using matched batch controls and optimized native
FFT libraries. In a scalar 4096-point comparison, distributing the same
implementation reduces the typical peak block time from
\StudyBatchPeak{} to \StudyCorePeak\,$\mu$s while increasing aggregate
instrumented cost by 18\%. With four aligned analyzers on one Rack engine
thread, the distributed implementation reduces the corresponding peak from
\StudyHostBatchPeak\,$\mu$s for native vDSP batch processing to
\StudyHostCorePeak\,$\mu$s, at the cost of more work and about 30\,ms of
additional publication delay. The measurements demonstrate a useful tradeoff
between processing bursts, total cost, and result latency; they do not show a
consistent reduction in missed deadlines.
\end{abstract}
